\section{Introduction}

For a prime \(p\) and an integer \(q\) not divisible by \(p\), let
\(\ord_p(q)\) be the multiplicative order of \(q\) in
\(\mathbb F_p^\times\).  Thus \(q\) is a primitive root modulo \(p\) when
\(\ord_p(q)=p-1\).  We study primes for which every member of a fixed finite
collection of positive prime bases is a primitive root.

The problem for one base is a case of Artin's primitive-root conjecture,
recorded by Hasse in 1927 \cite[Section~4.8]{HasseDiaries2012}.
Hooley proved the predicted asymptotic for one fixed integer under suitable
generalized Riemann hypotheses \cite{Hooley1967}.
Matthews obtained density results for simultaneous primitive roots under
generalized Riemann hypotheses \cite{Matthews1976}.  Anwar and Pappalardi
later established an infinitude criterion under Schinzel's Hypothesis~H
\cite{AnwarPappalardi2017}.
Unconditional results allowing a finite choice of bases were obtained by Gupta and Murty
\cite{GuptaMurty1984} and by Heath-Brown \cite{HeathBrown1986}.
In particular, Heath-Brown's theorem implies that at least one member of
any triple of distinct positive primes is a primitive root modulo
infinitely many primes.

Our result is an implication from four statements taken from the companion
manuscript \cite{PrimitiveRoot2}: a uniform Hecke zero-free assertion
(Input~\ref{input:hecke}), a marked Type~II estimate
(Input~\ref{input:marked}), a block sieve (Input~\ref{input:block}), and a
rough-number density estimate (Input~\ref{input:density}).  We reproduce
these statements in the precise forms used here.  They are hypotheses of
the following theorem; their proofs in \cite{PrimitiveRoot2} and its cited
companion manuscripts are outside the scope of this paper.

\begin{theorem}\label{thm:main}
Assume Inputs~\ref{input:hecke}, \ref{input:marked},
\ref{input:block}, and~\ref{input:density}.  Let \(k\ge1\) and let
\(q_1,\ldots,q_k\) be distinct positive primes.  There are constants
\(c_{\mathbf q}>0\) and \(x_{\mathbf q}>1\) such that, for every real
\(x\ge x_{\mathbf q}\), at least
\[
        c_{\mathbf q}\frac{x}{(\log x)^2}
\]
primes \(p\in(x,2x)\) satisfy
\[
        p\nmid q_1\cdots q_k,
        \qquad
        \ord_p(q_i)=p-1\quad(1\le i\le k).
\]
\end{theorem}

We extend the construction of primes with controlled predecessors in
\cite[Section~12]{PrimitiveRoot2} to a fixed progression.  The marked
Type~II estimate remains available after multiplication by a fixed
Dirichlet character, and Bombieri--Vinogradov supplies the required
distribution in that progression.  A direct sieve for the final balanced
semiprimes keeps its constant independent of the later marking parameters.
The progression makes every prescribed base a quadratic nonresidue.
Uniform splitting estimates in the individual Kummer fields then remove
the remaining prime divisors of each order index.  Section~2 gives this
reduction; the subsequent sections prove the two estimates it uses.
